%% otimiza_tx.tex
%% Otimização do Transmissor e Estratégia Sequencial Tx -> Rx
%% no SAR Biestático Helicoidal — Exploração do Ângulo de Brewster
%%
%% Idioma: Português
%% Compilar: pdflatex -interaction=nonstopmode otimiza_tx.tex  (duas passagens)
%% Figuras:  hypotheses/figures/otx_fig0X_*.pdf
%%           (gerar com utils/generate_figures_tx.py)

\documentclass[12pt,a4paper,oneside]{book}

%% ─── Pacotes ─────────────────────────────────────────────────────────────────
\usepackage[utf8]{inputenc}
\usepackage[T1]{fontenc}
\usepackage[portuguese]{babel}
\usepackage{lmodern}
\usepackage{microtype}
\usepackage{geometry}
\geometry{left=3cm, right=2.5cm, top=3cm, bottom=2.5cm}

\usepackage{amsmath,amssymb,amsfonts,amsthm}
\usepackage{mathtools}
\usepackage{bm}
\usepackage{siunitx}
\sisetup{output-decimal-marker={,}}

\usepackage{graphicx}
\graphicspath{{figures/}}

\usepackage{float}
\usepackage{caption}
\usepackage{subcaption}
\usepackage[table,xcdraw]{xcolor}
\usepackage{booktabs}
\usepackage{array}
\usepackage{multirow}
\usepackage{longtable}
\usepackage{tabularx}
\usepackage{setspace}
\onehalfspacing

\usepackage{algorithm}
\usepackage{algpseudocode}

\usepackage{tcolorbox}
\tcbuselibrary{theorems,skins,breakable}

\usepackage[unicode,colorlinks=true,
  linkcolor=blue!60!black,
  citecolor=green!50!black,
  urlcolor=cyan!60!black,
  pdftitle={Otimizacao do Transmissor e Estrategia Sequencial Tx-Rx em SAR Bistatico Helicoidal},
  pdfauthor={David Atencia}]{hyperref}

\usepackage{fancyhdr}
\setlength{\headheight}{15pt}
\pagestyle{fancy}
\fancyhf{}
\fancyhead[L]{\small\nouppercase{\leftmark}}
\fancyhead[R]{\small\thepage}
\renewcommand{\headrulewidth}{0.4pt}

%% ─── Caixas de destaque (mesmo estilo de explicacion_plan_de_vuelo.tex) ───────
\newcommand{\result}[2]{%
  \begin{tcolorbox}[enhanced, breakable,
    colback=yellow!8, colframe=orange!70!black,
    fonttitle=\bfseries,
    title={\protect#1},
    attach boxed title to top left={yshift=-2mm, xshift=5mm},
    boxed title style={colback=orange!70!black}]
    #2
  \end{tcolorbox}}

\newcommand{\nota}[2]{%
  \begin{tcolorbox}[enhanced, breakable,
    colback=green!5, colframe=green!50!black,
    fonttitle=\bfseries\small,
    title={\protect#1}]
    #2
  \end{tcolorbox}}

\newcommand{\atencion}[2]{%
  \begin{tcolorbox}[enhanced, breakable,
    colback=red!5, colframe=red!60!black,
    fonttitle=\bfseries\small,
    title={\protect#1}]
    #2
  \end{tcolorbox}}

%% ─── Teoremas ────────────────────────────────────────────────────────────────
\newtheorem{teorema}{Teorema}[chapter]
\newtheorem{proposicao}{Proposição}[chapter]
\newtheorem{lema}{Lema}[chapter]
\newtheorem{corolario}{Corolário}[chapter]
\theoremstyle{definition}
\newtheorem{definicaom}{Definição}[chapter]
\theoremstyle{remark}
\newtheorem*{observacao}{Observação}

%% ─── Macros ──────────────────────────────────────────────────────────────────
\newcommand{\RT}{R_T}
\newcommand{\RR}{R_R}
\newcommand{\thetaB}{\theta_B}
\newcommand{\nuno}{n_1}
\newcommand{\ndos}{n_2}
\newcommand{\Prcv}{P_r}
\newcommand{\Pt}{P_t}
\newcommand{\dxy}{\delta_{xy}}
\newcommand{\dz}{\delta_z}
\newcommand{\psiT}{\psi_{0,\mathrm{Tx}}}
\newcommand{\psiR}{\psi_{0,\mathrm{Rx}}}

% ─────────────────────────────────────────────────────────────────────────────
\begin{document}

\frontmatter

%%─── Capa ────────────────────────────────────────────────────────────────────
\begin{titlepage}
\centering
\vspace*{2cm}
{\LARGE\bfseries Otimização do Transmissor\\
no SAR Biestático Helicoidal:\\
uma Estratégia Sequencial Tx$\,\to\,$Rx\\
Guiada pelo Ângulo de Brewster}\\[1.2cm]
{\large\itshape Fundamentação física, formulação matemática e\\
sequência de desenvolvimento racional}\\[2cm]
{\normalsize Documento técnico complementar à dissertação\\
\textit{Otimização da Trajetória do Receptor em SAR Biestático\\
Helicoidal} e ao documento \texttt{explicacion\_plan\_de\_vuelo.tex}}\\[2cm]
{\normalsize David Atencia}\\[0.4cm]
{\normalsize 2026}
\vfill
\end{titlepage}

%%─── Resumo ──────────────────────────────────────────────────────────────────
\chapter*{Resumo}
\addcontentsline{toc}{chapter}{Resumo}

O sistema de otimização de trajetória descrito em \texttt{explicacion\_plan\_de\_vuelo.tex}
e no Capítulo~5 da dissertação otimiza \textbf{apenas} a posição do receptor (Rx) a
cada pulso, guiando-a pelo ângulo de Brewster $\thetaB$ para maximizar a
transmitância $T_2$ no trecho alvo$\to$Rx. O transmissor (Tx), no entanto, segue
uma trajetória helicoidal cônica cujos parâmetros geométricos
($\rho_0$, $z_0$, e portanto o ângulo de \textit{look} $\psi_0$) são fixados
\emph{a priori}, sem qualquer critério relacionado ao ângulo de Brewster — apesar
de a mesma equação de radar biestático depender, de forma inteiramente simétrica,
da transmitância $T_1$ no trecho Tx$\to$alvo.

Este documento demonstra que essa assimetria é uma lacuna de projeto, não uma
necessidade física, e propõe uma estratégia de \textbf{duas etapas}: primeiro
projetar o Tx (Etapa 1) para equilibrar a proximidade ao ângulo de Brewster
com a resolução vertical $\dz$ que sua própria abertura sintética determina;
depois, com o Tx fixado, otimizar o Rx (Etapa 2) exatamente como já é feito
hoje. Prova-se que a potência bistática recebida se fatora
\emph{exatamente} em um produto de um fator dependente apenas do Tx e um
fator dependente apenas do Rx — o que torna a decomposição sequencial
\emph{ótima, sem perda,} para o termo de potência, e uma aproximação bem
controlada para o termo de resolução. A metodologia é verificada
numericamente com os parâmetros reais do sistema já em uso
(\texttt{radarTx\_espiral\_plano\_voo.json}): o ângulo de \textit{look}
atualmente empregado, $\psi_0=58^\circ$, está a menos de $1^\circ$ do valor que a
otimização formal da Etapa 1 prescreve sob a mesma convenção de peso
logarítmico ($\gamma=0{,}3$) já usada para o Rx — evidência de que o projeto
atual já era \emph{quase} ótimo por construção, e de que este documento
formaliza e generaliza o critério que faltava.

%%─── Sumário ─────────────────────────────────────────────────────────────────
\tableofcontents

\mainmatter

%%═══════════════════════════════════════════════════════════════════════════════
\chapter{Introdução e Motivação}
%%═══════════════════════════════════════════════════════════════════════════════

\section{O problema: uma assimetria no projeto atual}
\label{sec:motivacao}

A equação de radar biestático com refração já foi derivada e está em uso no
sistema de otimização do receptor (documento \texttt{explicacion\_\allowbreak
plan\_\allowbreak de\_\allowbreak vuelo.tex}, Cap.~4, e função
\texttt{objective\_\allowbreak function.m}):

\begin{equation}
  \boxed{\Prcv = \frac{\Pt\,G_T\,G_R\,\sigma\,T_1\,T_2\,\lambda^2}{(4\pi)^3\,\RT^2\,\RR^2}}
  \label{eq:radar_bistatico}
\end{equation}

onde $T_1=T_1(\theta_{i,1})$ é a transmitância de potência TM no trecho
Tx$\to$alvo, avaliada no ângulo de incidência $\theta_{i,1}$ que o \emph{Tx}
forma com a normal à interface; $T_2=T_2(\theta_{i,2})$ é a transmitância
homóloga no trecho alvo$\to$Rx; e $\RT$, $\RR$ são os caminhos ópticos totais
de cada trecho (Tx$\to Q_1^*\to$alvo e alvo$\to Q_2^*\to$Rx,
respectivamente).

A Eq.~\eqref{eq:radar_bistatico} é \textbf{estruturalmente simétrica} em Tx e
Rx: trocar os papéis de $(T_1,\RT)$ e $(T_2,\RR)$ não altera a física. No
entanto, o sistema implementado hoje trata os dois lados de forma
completamente diferente:

\begin{itemize}
  \item \textbf{Rx}: a cada posição do Tx, a posição $(x,y)$ do Rx é
  escolhida por um otimizador não-convexo (\texttt{fmincon}, 5 candidatos de
  arranque) que maximiza uma função de custo multi-objetivo guiada
  explicitamente pelo ângulo de Brewster — o candidato de arranque \#4/5
  é construído a partir da distância de Brewster $d_B$
  (\texttt{buildRxSearchGeometry.m}, \texttt{chooseCandidateRxPos.m}).

  \item \textbf{Tx}: a trajetória (uma hélice cônica com raio menor/maior e
  altura menor/maior) é lida diretamente de um arquivo de parâmetros
  (\texttt{radarTx\_espiral\_plano\_voo.json}) e usada sem qualquer etapa de
  otimização. Nenhuma rotina do sistema calcula, verifica ou otimiza a
  proximidade entre o ângulo de incidência do Tx e $\thetaB$.
\end{itemize}

\nota{Evidência direta no código}{
O próprio \texttt{explicacion\_plan\_de\_vuelo.tex} (Cap.~1) descreve o Tx
como seguindo ``uma trajetória predeterminada \,--\, a hélice cônica \,--\,
otimizada para maximizar a diversidade angular e a resolução 3D''. Essa
frase é imprecisa: a hélice é escolhida para satisfazer requisitos de
resolução (via a condição $\beta=\psi_0$, Seção~\ref{sec:helix_review}),
mas \emph{não} há, em nenhum lugar do \texttt{system\_espiral\_plano\_voo.json}
ou do \texttt{run\_plano\_de\_voo.m}, um passo que escolha $\rho_0,z_0$ em
função de $\thetaB$.
}

\section{A pergunta central deste documento}

\begin{center}
\fbox{\parbox{0.85\textwidth}{\centering\itshape
Como projetar a trajetória do Tx para aproveitar o ângulo de Brewster tanto
quanto o Rx já aproveita — e por que faz sentido resolver primeiro o Tx e
só depois o Rx, em vez de otimizar os dois simultaneamente?
}}
\end{center}

Este documento responde a essa pergunta em três partes:

\begin{enumerate}
  \item Uma \textbf{fundamentação formal} (Capítulo~\ref{cap:fundamentacao})
  de por que a decomposição sequencial Tx$\to$Rx é justificada — mostrando
  que a potência bistática se fatora \emph{exatamente} em um termo
  só-Tx e um termo só-Rx, e que o único acoplamento residual (via a
  resolução horizontal $\dxy$) é assimilável pela liberdade posicional do
  Rx.

  \item Uma \textbf{formulação matemática} (Capítulos~\ref{cap:etapa1}
  e~\ref{cap:etapa2}) para cada etapa, reaproveitando integralmente o
  aparato já validado (coeficientes de Fresnel TM, fórmulas de resolução
  3D, função de custo logarítmica) e estendendo-o de forma simétrica ao
  Tx.

  \item Uma \textbf{verificação numérica} com os parâmetros reais do
  sistema, mostrando que a metodologia proposta reproduz — e generaliza —
  o valor de $\psi_0$ já em uso.
\end{enumerate}

\section{Relação com outros documentos}

Este documento pressupõe familiaridade com:
\begin{itemize}
  \item \texttt{explicacion\_resolucion.tex} — derivação das fórmulas de
  resolução 3D $\dxy,\dz$ a partir do $k$-espaço biestático.
  \item \texttt{explicacion\_plan\_de\_vuelo.tex} — derivação da equação de
  radar biestático com refração e da metodologia de otimização do Rx
  (função de custo, busca multi-início).
  \item \texttt{dissertacao\_tomo\_sar.tex} — Capítulos~2--3 (Fresnel TM,
  Brewster, resolução 3D) e Capítulo~5 (otimização do Rx), dos quais este
  documento herda notação e resultados já demonstrados, sem os
  redemonstrar.
\end{itemize}
Nenhum resultado já provado nesses documentos é redemonstrado aqui; eles
são citados e reutilizados como base para os resultados novos
(Capítulos~\ref{cap:fundamentacao}--\ref{cap:etapa1}).

%%═══════════════════════════════════════════════════════════════════════════════
\chapter{Fundamentos Reaproveitados}
\label{cap:fundamentos}
%%═══════════════════════════════════════════════════════════════════════════════

Este capítulo reúne, sem redemonstrar, os resultados da dissertação e dos
documentos irmãos que serão usados como blocos de construção.

\section{Ângulo de Brewster e transmitância TM}

Para incidência TM em meios sem perdas e não-magnéticos, o único ângulo de
incidência $\theta_i$ que anula a reflectância é
\begin{equation}
  \thetaB = \arctan\!\left(\frac{\ndos}{\nuno}\right)
  \label{eq:brewster}
\end{equation}
(Teorema~2.1 da dissertação). A transmitância de potência correspondente é
\begin{equation}
  T(\theta_i) = \frac{\ndos\cos\theta_t}{\nuno\cos\theta_i}\,|t_{TM}(\theta_i)|^2
  = 1-|r_{TM}(\theta_i)|^2,
  \label{eq:fresnel_T}
\end{equation}
com $T(\thetaB)=1$. Para $\nuno=1,\ndos=2$ (solo com $\varepsilon_r=4$),
$\thetaB\approx 63{,}43^\circ$. A função $T(\theta_i)$ é \textbf{unimodal}: cresce
monotonicamente de $\theta_i=0$ até $\thetaB$ e decresce monotonicamente daí
até $\theta_i=90^\circ$ (visível na Fig.~2.1 da dissertação) — propriedade que
será usada nos Capítulos~\ref{cap:fundamentacao}--\ref{cap:etapa1}.

\section{Geometria da hélice cônica}

Cada antena (Tx ou Rx) descreve, na formulação idealizada do documento
de resolução espacial e do Capítulo~3 da dissertação, uma hélice cônica
caracterizada por:
\begin{itemize}
  \item $\rho_0$ — raio médio (entre \texttt{RaioMenorEsp} e
  \texttt{RaioMaiorEsp});
  \item $z_0$ — altura média (entre \texttt{AltMenorEsp} e
  \texttt{AltMaiorEsp});
  \item $\psi_0=\arctan(\rho_0/z_0)$ — ângulo de \textit{look} médio;
  \item $R_0=\sqrt{\rho_0^2+z_0^2}$ — alcance oblíquo médio;
  \item $\beta=\arctan(\Delta z/\Delta\rho)$ — ângulo de inclinação da
  hélice, com $\Delta\rho,\Delta z$ as excursões de raio e altura;
  \item $B_{helix}=\sqrt{\Delta\rho^2+\Delta z^2}$ — comprimento de arco da
  hélice (abertura sintética total).
\end{itemize}

\section{Condição ótima de inclinação (revisão)}
\label{sec:helix_review}

\begin{proposicao}[já demonstrada — dissertação, Seção~3.6]
A projeção perpendicular da hélice sobre a linha de visada,
$B_\perp = B_{helix}\,|\cos(\beta-\psi_0)|$, é maximizada — e portanto $\dz$
é minimizado, para $B_{helix}$ fixo — quando $\beta=\psi_0$, resultando em
$B_\perp=B_{helix}$.
\end{proposicao}

Este documento \textbf{não questiona} essa condição; ela é adotada como
válida e usada throughout (Capítulo~\ref{cap:etapa1}) como restrição de
igualdade $\beta=\psi_0$ ao dimensionar o Tx.

\section{Fórmulas de resolução 3D (revisão)}
\label{sec:resolucao_review}

Da Seção~3.4--3.5 da dissertação (Eqs.~3.10--3.11) e de
\texttt{calculateBistaticResolution.m}:
\begin{align}
  \dz &= \frac{c}{2\,W_z}, \qquad
  W_z = \ndos B\cos\theta_{t,0}
  + \frac{c\,B_\perp\sin\psi_0\cos\psi_0}{\lambda_0 R_0 \ndos\cos\theta_{t,0}}
  \label{eq:Wz_review}\\
  \dxy(\Delta\phi) &= \frac{0{,}60\,\lambda_0}{\pi\,\sin\psi_0\,|\cos(\Delta\phi/2)|}
  \label{eq:dxy_review}
\end{align}
com $\cos\theta_{t,0}=\sqrt{1-\sin^2\psi_0/\ndos^2}$ (Snell), $B$ a largura
de banda do chirp e $\Delta\phi$ o NorthOffset (separação azimutal Tx--Rx
vista do alvo).

\begin{teorema}[já demonstrado — dissertação, Seção~3.4]
\label{teo:invariancia_review}
$\dz$ é invariante a $\Delta\phi$: a Eq.~\eqref{eq:Wz_review} não contém
$\Delta\phi$.
\end{teorema}

Este teorema é a pedra angular do Capítulo~\ref{cap:fundamentacao}: ele
garante que a resolução vertical não depende de \emph{onde}, azimutalmente,
o Rx está em relação ao Tx — apenas do módulo do ângulo de \textit{look}
$\psi_0$ envolvido.

\section{A arquitetura computacional já implementada}
\label{sec:arquitetura_review}

O código de produção já separa, estruturalmente, o cálculo do trecho
Tx$\to$alvo do trecho alvo$\to$Rx. Em \texttt{precomputeTxData.m}:

\begin{quote}\ttfamily\footnotesize
\% PRECOMPUTETXDATA\ \ Pre-computes Tx→Target refraction\\
\% data for a fixed Tx position.\\
\%\\
\% The Tx→Target leg is independent of the Rx position\\
\% being optimized by fmincon. Pre-computing it once per\\
\% Tx position avoids recomputing the Fermat refraction\\
\% point for every objective\_function evaluation ---\\
\% eliminating \textasciitilde 50\% of Fermat calls.
\end{quote}

Ou seja: $T_1$ e $\RT$ já são calculados \textbf{uma única vez por posição
de Tx}, \emph{antes} do laço de otimização do Rx, e reutilizados como
constantes em toda avaliação de \texttt{objective\_function} durante essa
otimização (linhas 43--45 de \texttt{objective\_function.m}). O
Capítulo~\ref{cap:fundamentacao} formaliza matematicamente por que essa
escolha de arquitetura é não apenas conveniente, mas \emph{correta}.

%%═══════════════════════════════════════════════════════════════════════════════
\chapter{Fundamentação da Sequência Tx$\,\to\,$Rx}
\label{cap:fundamentacao}
%%═══════════════════════════════════════════════════════════════════════════════

\section{Fatoração exata da potência bistática}

\begin{proposicao}[Separabilidade multiplicativa de $\Prcv$]
\label{prop:separabilidade}
Para um alvo $P$ fixo, com a posição do Tx e a posição do Rx como variáveis
independentes, a potência bistática~\eqref{eq:radar_bistatico} se fatora
como
\begin{equation}
  \Prcv(\mathrm{Tx},\mathrm{Rx};P) = \underbrace{\left[\frac{\Pt\,\sigma\,\lambda^2}{(4\pi)^3}\right]}_{\text{constante}}
  \cdot\underbrace{\frac{G_T(\mathrm{Tx})\,T_1(\mathrm{Tx})}{\RT(\mathrm{Tx})^2}}_{f_{\mathrm{Tx}}(\mathrm{Tx})}
  \cdot\underbrace{\frac{G_R(\mathrm{Rx})\,T_2(\mathrm{Rx})}{\RR(\mathrm{Rx})^2}}_{f_{\mathrm{Rx}}(\mathrm{Rx})}
  \label{eq:fatoracao}
\end{equation}
onde $f_{\mathrm{Tx}}$ depende \emph{apenas} da posição do Tx (e de $P$), e
$f_{\mathrm{Rx}}$ depende \emph{apenas} da posição do Rx (e de $P$).
\end{proposicao}

\begin{proof}
Por construção geométrica (\texttt{explicacion\_plan\_de\_vuelo.tex},
Cap.~4): $\RT=R_1+R_2$ e $\theta_{i,1}$ (logo $T_1$) são determinados
inteiramente pelo segmento Tx$\to Q_1^*\to P$, que não envolve a posição
do Rx. Analogamente, $\RR=R_3+R_4$ e $\theta_{i,2}$ (logo $T_2$) são
determinados inteiramente pelo segmento $P\to Q_2^*\to$Rx, que não envolve
a posição do Tx. O ganho de antena $G_T$ é o padrão de radiação do Tx
avaliado no ângulo Tx$\to P$ (\texttt{calculateTxGainsForTargets}), e
$G_R$ o padrão do Rx avaliado no ângulo $P\to$Rx
(\texttt{getAntennaGain}) — nenhum dos dois mistura as duas geometrias.
Substituindo em~\eqref{eq:radar_bistatico} e agrupando os fatores por
dependência, obtém-se~\eqref{eq:fatoracao}.
\end{proof}

\begin{corolario}
\label{cor:otimizacao_separavel}
Para o termo de potência pura ($\alpha_{\mathrm{res}}=0$ na função de
custo), maximizar $\Prcv$ conjuntamente sobre $(\mathrm{Tx},\mathrm{Rx})$
é \textbf{exatamente equivalente} a maximizar $f_{\mathrm{Tx}}$ e
$f_{\mathrm{Rx}}$ separadamente:
\[
  \operatorname*{arg\,max}_{\mathrm{Tx},\mathrm{Rx}} \Prcv
  = \Big(\operatorname*{arg\,max}_{\mathrm{Tx}} f_{\mathrm{Tx}},\ \operatorname*{arg\,max}_{\mathrm{Rx}} f_{\mathrm{Rx}}\Big),
\]
\emph{sem nenhuma perda}, independentemente da ordem em que as duas
otimizações são feitas.
\end{corolario}

Este é o resultado mais forte deste documento: para o objetivo de potência,
a decomposição Tx$\to$Rx (ou Rx$\to$Tx — a ordem, isoladamente, é
indiferente para este termo) não é uma aproximação, é uma identidade
algébrica.

\section{O acoplamento residual: resolução}
\label{sec:acoplamento}

A função de custo completa (Cap.~5 da dissertação;
\texttt{objective\_function.m}) inclui o termo de resolução:
\begin{equation}
  J = -(1-\alpha_{\mathrm{res}})\log_{10}\Prcv + \alpha_{\mathrm{res}}\log_{10}(\dxy\cdot\dz).
  \label{eq:custo_geral}
\end{equation}
Diferentemente de $\Prcv$, o termo $\dxy\cdot\dz$ \textbf{não} se fatora de
forma limpa: pela Eq.~\eqref{eq:Wz_review}, $\dz$ depende de $\psi_0$ e de
$(\beta,B_{helix})$; pela Eq.~\eqref{eq:dxy_review}, $\dxy$ depende de
$\psi_0$ \emph{e} de $\Delta\phi$ (que mistura os azimutes de Tx e Rx).
Três observações controlam o impacto desse acoplamento:

\begin{enumerate}
  \item \textbf{$\dz$ não depende de $\Delta\phi$}
  (Teorema~\ref{teo:invariancia_review}) — o único acoplamento
  Tx$\leftrightarrow$Rx nele é através de \emph{qual} $\psi_0$ se usa para
  avaliar a Eq.~\eqref{eq:Wz_review}.

  \item \textbf{$(\beta,B_{helix})$ são fixados pelo Tx antes de qualquer
  otimização do Rx} — literalmente, no código de produção,
  \texttt{B\_helix} e \texttt{beta\_helix} são calculados uma vez
  (\texttt{run\_plano\_de\_voo.m}, linhas 61--65) a partir apenas dos
  parâmetros geométricos do Tx, \emph{antes} do laço \texttt{for tx\_idx}
  que otimiza o Rx, e são passados como constantes para
  \texttt{objective\_function} em toda iteração subsequente.

  \item \textbf{$\dxy$ retém liberdade suficiente no Rx} — dado
  $(\beta,B_{helix})$ fixos pelo Tx, o Rx ainda controla livremente
  $\Delta\phi$ (via sua posição azimutal) e sua própria contribuição a
  $\psi_0$ na Eq.~\eqref{eq:dxy_review}, exatamente como já demonstrado
  pela otimização de duas dimensões ($x,y$) já em produção.
\end{enumerate}

\begin{proposicao}[Decomposição hierárquica]
\label{prop:decomposicao_hierarquica}
Definindo a Etapa 1 como a escolha das constantes de projeto do Tx
$(\rho_0,z_0)\Rightarrow(\psi_0,\beta{=}\psi_0,R_0,B_{helix})$ e a Etapa 2
como a otimização usual do $(x,y)$ do Rx \emph{dado} essas constantes, a
sequência Etapa~1$\to$Etapa~2:
\begin{itemize}
  \item é \textbf{exatamente ótima} para o termo de potência
  (Corolário~\ref{cor:otimizacao_separavel});
  \item é \textbf{exatamente consistente} com a invariância de $\dz$ a
  $\Delta\phi$ (Teorema~\ref{teo:invariancia_review}) — a Etapa 1 não
  precisa se preocupar com $\Delta\phi$, que é resolvido inteiramente na
  Etapa 2;
  \item deixa para a Etapa 2 exatamente as variáveis que ela já otimiza
  hoje ($x,y$ do Rx, logo $\Delta\phi$ e a contribuição do Rx a
  $\dxy$), sem removida nenhuma liberdade que o otimizador atual utilize.
\end{itemize}
\end{proposicao}

\begin{observacao}
A única aproximação genuína introduzida é que a Etapa 1 dimensiona
$\dz$ usando o $\psi_0$ \emph{nominal do Tx} ($\psiT$), enquanto a Etapa 2
(como já implementado em \texttt{calculateBistaticResolution.m}) reavalia
$\dz,\dxy$ usando o $\psi_0$ \emph{realizado pelo Rx} ($\psiR$) — os dois
podem diferir porque a posição ótima do Rx não é, em geral, uma réplica
angularmente OFF-SET da hélice do Tx. Essa é precisamente a mesma
aproximação de "ângulo de \textit{look}'' que a dissertação já documenta e
quantifica em $\sim\!7\text{--}8\%$ de erro (§6.4.2, análise de erros de
$\dxy$ no eixo). O impacto é, portanto, do mesmo tipo e ordem de grandeza
de uma aproximação já aceita e caracterizada em outro lugar do corpus —
não um erro novo introduzido por esta proposta.
\end{observacao}

\section{Por que esta ordem, e não a ordem inversa ou a otimização conjunta}

\begin{itemize}
  \item \textbf{Tx antes de Rx, e não o contrário}: $(\beta,B_{helix})$
  são propriedades físicas da trajetória do Tx que devem ser fixadas
  \emph{antes} do voo (a hélice é executada mecanicamente pela
  plataforma); a posição do Rx, em contraste, pode ser recalculada
  pulso-a-pulso. Logo, na prática de engenharia, o Tx \emph{tem} que ser
  decidido primeiro — este documento apenas formaliza \emph{como}
  decidi-lo racionalmente, em vez de por herança arbitrária de um arquivo
  de parâmetros.

  \item \textbf{Sequencial, e não conjunta}: pela
  Proposição~\ref{prop:separabilidade}, uma otimização conjunta
  $(\mathrm{Tx},\mathrm{Rx})\in\mathbb{R}^2\times\mathbb{R}^2$ não
  encontraria nenhum ótimo de potência melhor do que a decomposição
  sequencial — apenas pagaria um custo computacional maior (busca em
  4 dimensões com múltiplos mínimos locais em vez de duas buscas
  independentes de dimensão menor, uma delas escalar; ver
  Capítulo~\ref{cap:algoritmo}). A única situação em que a otimização
  conjunta poderia, em princípio, superar a sequencial é através do termo
  de resolução residual da Seção~\ref{sec:acoplamento} — cujo impacto foi
  limitado, acima, à mesma ordem de grandeza de uma aproximação já aceita
  no corpus.
\end{itemize}

%%═══════════════════════════════════════════════════════════════════════════════
\chapter{Etapa 1 --- Otimização do Transmissor}
\label{cap:etapa1}
%%═══════════════════════════════════════════════════════════════════════════════

\section{Variáveis de projeto}

O projetista escolhe $\rho_0,z_0$ (equivalentemente $\psiT=\arctan(\rho_0/z_0)$
e $R_0=\sqrt{\rho_0^2+z_0^2}$) e o comprimento de hélice $B_{helix}$
(determinado pela excursão $\Delta\rho,\Delta z$ desejada para a apertura
sintética). A condição $\beta=\psiT$ (Seção~\ref{sec:helix_review}) é
imposta por construção. Resta, portanto, um problema de otimização
essencialmente \textbf{escalar} em $\psiT$ (para $R_0$ ou $z_0$ fixado por
uma restrição de missão — Seção~\ref{sec:restricoes}) — muito mais simples
do que o problema não-convexo de duas dimensões que a Etapa 2 já resolve.

\section{O conflito fundamental: $45^\circ$ contra $\thetaB$}
\label{sec:conflito}

O ângulo de incidência do Tx no alvo, $\theta_{i,1}$, coincide com $\psiT$
em primeira ordem, para um alvo próximo ao centro da malha (a mesma
aproximação de \textit{look angle} já usada em toda a análise de resolução
— Observação da Seção~\ref{sec:acoplamento}). Logo, $T_1(\psiT)$ é
maximizado quando
\begin{equation}
  \psiT = \thetaB.
  \label{eq:otimo_potencia}
\end{equation}

Por outro lado, o termo tomográfico de $W_z$ na
Eq.~\eqref{eq:Wz_review} contém o fator $\sin\psi_0\cos\psi_0
=\tfrac12\sin(2\psi_0)$, maximizado exatamente em
\begin{equation}
  \psi_0 = 45^\circ.
  \label{eq:otimo_tomografico}
\end{equation}
Além disso, o termo de largura de banda $\ndos B\cos\theta_{t,0}$
\emph{decresce monotonicamente} com $\psi_0$ (pois
$\cos\theta_{t,0}=\sqrt{1-\sin^2\psi_0/\ndos^2}$ é decrescente em
$\psi_0\in(0^\circ,90^\circ)$), reforçando ainda mais o ótimo tomográfico abaixo de
$45^\circ$-$63^\circ$.

\begin{teorema}[Conflito estrutural]
\label{teo:conflito}
Para $\ndos>\nuno$ (interface ar-solo típica), $\thetaB=\arctan(\ndos/\nuno)>45^\circ$
sempre que $\ndos>\nuno$. Logo $\thetaB\ne 45^\circ$ e os ótimos individuais de
$T_1$ e de $W_z$ (via seu termo tomográfico) \textbf{nunca coincidem}
exatamente, para nenhum par de meios fisicamente distintos.
\end{teorema}

\begin{proof}
$\arctan(x)=45^\circ\Leftrightarrow x=1\Leftrightarrow \ndos=\nuno$, excluído por
hipótese ($\ndos\ne\nuno$, caso contrário não há interface física).
\end{proof}

Este é o análogo, no lado do Tx, do trade-off potência--resolução que
motiva a função de custo do Rx (Cap.~5.4 da dissertação); o
Teorema~\ref{teo:conflito} mostra que ele é \textbf{inevitável}, não uma
peculiaridade de parâmetros específicos.

\subsection{Verificação numérica}

A Figura~\ref{fig:psi0_tradeoff} mostra $T_1(\psi_0)$ e $\dz(\psi_0)$ para
os parâmetros reais do sistema
($f_0=400\,\mathrm{MHz}$, $B=50\,\mathrm{MHz}$, $\ndos=2$,
$B_{helix}=47{,}17\,\mathrm{m}$), sob duas restrições de projeto distintas:
alcance oblíquo $R_0$ fixo (relevante quando o orçamento de enlace limita
o alcance máximo) e altura de voo $z_0$ fixa (relevante quando um teto
operacional limita a altitude).

\begin{figure}[H]
\centering
\includegraphics[width=0.98\textwidth]{otx_fig01_psi0_tradeoff.pdf}
\caption{Conflito $\psi_0=45^\circ$ (ótimo tomográfico, minimiza $\dz$) vs.\
$\psi_0=\thetaB=63{,}4^\circ$ (ótimo de Brewster, $T_1=1$), calculado com os
parâmetros reais do sistema. Em ambos os cenários de restrição, o
ângulo atualmente usado ($\psi_0=58^\circ$) já se situa \emph{dentro} do
platô largo onde $T_1>0{,}99$.}
\label{fig:psi0_tradeoff}
\end{figure}

Uma propriedade notável, visível na Fig.~\ref{fig:psi0_tradeoff}: $T_1$ é
\textbf{larga} perto de $\thetaB$ (para $\ndos/\nuno=2$, $T_1\ge 0{,}90$ para
$\psi_0\in[18^\circ,77^\circ]$), enquanto o pico tomográfico de $W_z$ é mais estreito.
Isso significa que o custo, em potência, de priorizar a resolução é
tipicamente pequeno; o Capítulo formaliza esse trade-off para escolher o
ponto de equilíbrio correto em vez de confiar nessa observação
qualitativamente.

\section{Função de custo do Tx}

Em analogia direta com a função de custo do Rx (Eq.~\eqref{eq:custo_geral}),
define-se
\begin{equation}
  \boxed{J_{\mathrm{Tx}}(\psiT;\gamma) = -(1-\gamma)\log_{10} T_1(\psiT)
  + \gamma\,\log_{10}\dz(\psiT)}
  \label{eq:custo_tx}
\end{equation}
com $\gamma\in[0,1]$ o peso de resolução do \textbf{lado do Tx}
(analogamente a $\alpha_{\mathrm{res}}$ para o Rx, mas um parâmetro
distinto — nada obriga $\gamma=\alpha_{\mathrm{res}}$, embora a
Seção~\ref{sec:validacao} mostre que adotar o mesmo valor numérico produz
um resultado consistente com o sistema já em uso). A normalização
logarítmica é a mesma já justificada no Cap.~5 da dissertação: $T_1\in(0,1]$
e $\dz\sim\!10^2\,\mathrm{cm}$ têm escalas e unidades incomensuráveis, e o
logaritmo as torna comparáveis.

\begin{proposicao}[Unimodalidade numérica]
Para os parâmetros do sistema (Seção~4.2.1) e $\gamma\in[0,1]$, $J_{\mathrm{Tx}}$
tem um único mínimo interior em $\psiT\in(0^\circ,90^\circ)$ — verificado por
varredura exaustiva (passo $0{,}04^\circ$) para $\gamma\in\{0{,}0.1,\dots,1{,}0\}$
(Fig.~\ref{fig:gamma_sweep}a). Isso decorre de $T_1(\psi_0)$ ser unimodal
com pico em $\thetaB$ (Seção~2.1) e $W_z(\psi_0)$ ser unimodal com pico
verificado numericamente entre $33^\circ$ e $45^\circ$, dependendo da restrição de
projeto (Seção~\ref{sec:restricoes}) — a combinação convexa (em
$\log_{10}$) de duas funções unimodais com um único vale cada, cujos vales
estão a menos de $30^\circ$ de distância, não introduz mínimos locais
espúrios neste intervalo.
\end{proposicao}

\begin{figure}[H]
\centering
\includegraphics[width=0.98\textwidth]{otx_fig02_gamma_sweep.pdf}
\caption{(a) $J_{\mathrm{Tx}}(\psiT;\gamma)$ para vários $\gamma$
(cenário $R_0$ fixo) — mínimo único em cada curva. (b) Ângulo ótimo
$\psiT^*(\gamma)$: decresce suavemente de $\thetaB$ ($\gamma=0$) até o
ótimo tomográfico ($\gamma=1$). O losango marca $\gamma=0{,}3$ — mesmo peso
já usado para $\alpha_{\mathrm{res}}$ no Rx — produzindo
$\psiT^*=57{,}6^\circ$, a $0{,}4^\circ$ do valor atualmente empregado ($58^\circ$).}
\label{fig:gamma_sweep}
\end{figure}

\section{Restrições práticas}
\label{sec:restricoes}

A Eq.~\eqref{eq:custo_tx} é minimizada sujeita a restrições de missão que
fixam a escala do problema (o par $(\rho_0,z_0)$ tem dois graus de
liberdade; $\psiT$ é só a razão entre eles):

\begin{itemize}
  \item \textbf{Alcance oblíquo fixo ($R_0=\text{const.}$)} — apropriado
  quando o orçamento de enlace (Seção~4.6, \texttt{explicacion\_plan\_de\_vuelo.tex})
  fixa o alcance máximo tolerável para a relação sinal-ruído desejada. Sob
  esta restrição, $\rho_0=R_0\sin\psiT$, $z_0=R_0\cos\psiT$.

  \item \textbf{Altitude de voo fixa ($z_0=\text{const.}$)} — apropriado
  quando um teto operacional (regulatório ou da plataforma) limita a
  altitude. Sob esta restrição, $\rho_0=z_0\tan\psiT$, $R_0=z_0/\cos\psiT$.

  \item \textbf{Cobertura da malha de alvos} — $\rho_0$ deve ser
  suficientemente grande para que a projeção horizontal da hélice
  circunde toda a malha de alvos (raio $\gtrsim$ metade da diagonal da
  malha), sob pena de perda de diversidade angular nas bordas.

  \item \textbf{Limites de elevação da antena} — $\psiT$ deve permanecer
  dentro da abertura de elevação do padrão de radiação
  (\texttt{AperturaElev}, tipicamente $\pm 25^\circ$ em torno do
  apontamento nominal), sob pena de $G_T$ cair fora do lóbulo principal.
\end{itemize}

Ambos os cenários numéricos (Fig.~\ref{fig:psi0_tradeoff}) são
apresentados lado a lado justamente porque a escolha entre eles depende de
qual restrição domina a missão específica — não há uma resposta universal,
e apresentar apenas o cenário mais favorável seria enganoso.

\section{Algoritmo de busca}

Ao contrário do problema do Rx (não-convexo, 2 dimensões, múltiplos
mínimos locais reais — Cap.~5.5 da dissertação), o problema da
Eq.~\eqref{eq:custo_tx} é escalar e unimodal (verificado acima), portanto
resolúvel por busca de seção áurea (\textit{golden-section search}) em
$O(\log(1/\varepsilon))$ avaliações, sem necessidade de multi-início:

\begin{algorithm}[H]
\caption{Otimização do Tx (Etapa 1)}
\begin{algorithmic}[1]
\Require $n_1,n_2$, $B$, $B_{helix}$, restrição de projeto ($R_0$ ou $z_0$),
  $\gamma$, limites $[\psi_{\min},\psi_{\max}]$
\State $\thetaB \gets \arctan(n_2/n_1)$
\State $\psiT^* \gets$ \Call{BuscaSecaoAurea}{$J_{\mathrm{Tx}}(\cdot;\gamma)$, $\psi_{\min}$, $\psi_{\max}$}
\State $\beta \gets \psiT^*$ \Comment{condição ótima já demonstrada (Seção~2.3)}
\State $(\rho_0,z_0) \gets$ resolver a partir de $\psiT^*$ e da restrição de projeto
\State \Return $\rho_0,z_0,\beta,B_{helix}$ \Comment{constantes para a Etapa 2}
\end{algorithmic}
\end{algorithm}

\section{Extensão a uma malha de alvos}

Para uma malha extensa (não um único alvo pontual), $\theta_{i,1}$ varia
por alvo. Em analogia direta ao estudo de sensibilidade já realizado para
o Rx (§5.8 da dissertação, grades de $30/60/100\,\mathrm{m}$), recomenda-se
avaliar $T_1$ no centroide da malha para o dimensionamento de $\psiT$
(Eq.~\eqref{eq:custo_tx}), e verificar \emph{a posteriori} a dispersão de
$T_1$ nas bordas da malha — exatamente como a dissertação já verifica a
dispersão de potência entre grades de tamanhos distintos. Este documento
não repete esse estudo numérico; ele indica que a mesma metodologia se
aplica sem modificação ao lado do Tx.

\section{Validação com os parâmetros reais do sistema}
\label{sec:validacao}

\begin{table}[H]
\centering
\caption{Etapa 1 aplicada aos parâmetros de
\texttt{radarTx\_espiral\_plano\_voo.json} (cenário $R_0$ fixo,
$R_0=188{,}7\,\mathrm{m}$), com $\gamma=0{,}3$.}
\label{tab:validacao_etapa1}
\begin{tabularx}{\textwidth}{@{}l r X@{}}
\toprule
Grandeza & Valor & Origem \\
\midrule
$\thetaB$ & $63{,}43^\circ$ & $\arctan(n_2/n_1)$, $n_2=2$ \\
$\psi_0$ atualmente em uso & $58{,}0^\circ$ & JSON de parâmetros do Tx (raio e altura médios da hélice) \\
$T_1$ no valor atual & $0{,}9938$ ($-0{,}027\,\mathrm{dB}$) & Eq.~\eqref{eq:fresnel_T} \\
$\psiT^*$ (Etapa 1, $\gamma=0{,}3$) & $57{,}6^\circ$ & Eq.~\eqref{eq:custo_tx}, busca de seção áurea \\
Diferença & $0{,}4^\circ$ & --- \\
\bottomrule
\end{tabularx}
\end{table}

\result{Interpretação}{
O valor de $\psi_0=58^\circ$ já usado no sistema não foi escolhido por um
critério de Brewster explícito — é herdado apenas da condição de resolução
$\beta=\psi_0$ combinada com uma escolha razoável, mas não formalizada, de
$\rho_0,z_0$. Ainda assim, sob o mesmo peso $\gamma=\alpha_{\mathrm{res}}=0{,}3$
já adotado no lado do Rx, a Etapa~1 formal prescreve $57{,}6^\circ$ — uma
diferença de $0{,}4^\circ$, dentro da margem de arredondamento dos parâmetros de
projeto. Isto \textbf{não} significa que a formalização seja
desnecessária: significa que o sistema já operava, por boa sorte de
projeto, próximo do ótimo formal para este conjunto específico de
parâmetros — e que qualquer mudança futura de $n_2$, $B$, $B_{helix}$ ou da
restrição de missão (Seção~\ref{sec:restricoes}) exigirá recalcular
$\psiT^*$, o que este documento agora torna possível fazer de forma
sistemática, em vez de por tentativa e erro.
}

Como contraponto, a Tabela~\ref{tab:validacao_pior_caso} mostra que a
formalização deixa de ser irrelevante assim que o ponto de partida se
afasta do platô de Brewster.

\begin{table}[H]
\centering
\caption{Ganho de potência ao corrigir escolhas de $\psi_0$ mais distantes
de $\thetaB$ — ganho obtido \emph{apenas} pela Etapa 1, sem custo algum
para a liberdade da Etapa 2 (Corolário~\ref{cor:otimizacao_separavel}).}
\label{tab:validacao_pior_caso}
\begin{tabular}{lrr}
\toprule
$\psi_0$ de partida & $T_1$ & Ganho ao mover para $\psiT^*=57{,}6^\circ$ \\
\midrule
$20^\circ$ & $0{,}9026$ & $+0{,}42\,\mathrm{dB}$ \\
$80^\circ$ & $0{,}8155$ & $+0{,}86\,\mathrm{dB}$ \\
\bottomrule
\end{tabular}
\end{table}

%%═══════════════════════════════════════════════════════════════════════════════
\chapter{Etapa 2 --- Otimização do Receptor Dado o Tx Ótimo}
\label{cap:etapa2}
%%═══════════════════════════════════════════════════════════════════════════════

\section{O que a Etapa 1 entrega à Etapa 2}

Ao final da Etapa 1, ficam fixados: $\rho_0,z_0$ (logo $\psiT$, $R_0$),
$\beta=\psiT$ e $B_{helix}$ — os mesmos quatro parâmetros que
\texttt{run\_plano\_de\_voo.m} já calcula uma única vez, antes do laço de
otimização do Rx (Seção~\ref{sec:arquitetura_review}). A diferença é que
esses valores deixam de ser lidos passivamente de um JSON e passam a ser o
resultado de uma otimização escalar explícita e justificada.

\section{A Etapa 2 é a metodologia já existente, sem alterações}

Dado o Tx fixo, o problema resolvido pulso a pulso é exatamente o já
descrito em \texttt{explicacion\_plan\_de\_vuelo.tex} e no Cap.~5 da
dissertação:
\begin{equation}
  \operatorname*{arg\,min}_{(x,y)_{\mathrm{Rx}}}
  \Big[-(1-\alpha_{\mathrm{res}})\log_{10}\Prcv + \alpha_{\mathrm{res}}\log_{10}(\dxy\cdot\dz)\Big],
  \label{eq:etapa2}
\end{equation}
resolvido por \texttt{fmincon} com 5 candidatos de arranque guiados
geometricamente (posição oposta, perpendiculares, região de Brewster).
Este documento não modifica essa etapa — ele apenas garante que, quando
ela roda, $T_1$ (fixo, pré-computado) já está o mais próximo possível de
$1$ que a Eq.~\eqref{eq:custo_tx} permite, e que $(\beta,B_{helix})$ já
refletem uma escolha racional, não arbitrária.

\section{Ganho aditivo: prova formal}

\begin{proposicao}[Aditividade do ganho de potência]
\label{prop:aditividade}
Sejam $\psiT^{(0)}$ o ângulo de \textit{look} do Tx antes da otimização
da Etapa~1 e $\psiT^{(1)}=\psiT^*$ o ângulo após. Para \emph{qualquer}
posição de Rx escolhida pela Etapa~2 (ótima ou não), o ganho de potência
recebida ao substituir $\psiT^{(0)}$ por $\psiT^{(1)}$ é
\begin{equation}
  \Delta \Prcv[\mathrm{dB}] = 10\log_{10}\!\left(\frac{T_1(\psiT^{(1)})}{T_1(\psiT^{(0)})}\right),
  \label{eq:ganho_aditivo}
\end{equation}
\textbf{independente} da posição do Rx escolhida.
\end{proposicao}

\begin{proof}
Pela Proposição~\ref{prop:separabilidade}, $\Prcv=K\cdot f_{\mathrm{Tx}}(\mathrm{Tx})\cdot f_{\mathrm{Rx}}(\mathrm{Rx})$
com $K$ constante. Para dois valores de $\psiT$ e a \emph{mesma} posição
de Rx, a razão $\Prcv^{(1)}/\Prcv^{(0)} = f_{\mathrm{Tx}}(\mathrm{Tx}^{(1)})/f_{\mathrm{Tx}}(\mathrm{Tx}^{(0)})$
não depende de $f_{\mathrm{Rx}}(\mathrm{Rx})$, que se cancela no
quociente. Como $\RT,G_T$ variam pouco entre $\psiT^{(0)},\psiT^{(1)}$
comparados a $T_1$ perto de $\thetaB$ (a componente dominante da variação
é a de $T_1$, que passa por um zero de reflectância), a razão é
$\approx T_1(\psiT^{(1)})/T_1(\psiT^{(0)})$, dando~\eqref{eq:ganho_aditivo}
em dB.
\end{proof}

\nota{Consequência prática}{
A Tabela~\ref{tab:validacao_pior_caso} não é apenas uma comparação
teórica: pela Proposição~\ref{prop:aditividade}, \emph{qualquer} ganho
ali listado se soma diretamente ao orçamento de enlace já estimado em
\texttt{explicacion\_plan\_de\_vuelo.tex} (Cap.~4, $\approx\!-32\,\mathrm{dBm}$
sob a hipótese $T_1=T_2=1$), sem exigir nenhuma mudança na Etapa~2 nem
recalibração de $\alpha_{\mathrm{res}}$.
}

%%═══════════════════════════════════════════════════════════════════════════════
\chapter{Algoritmo Consolidado}
\label{cap:algoritmo}
%%═══════════════════════════════════════════════════════════════════════════════

\section{Sequência completa}

\begin{algorithm}[H]
\caption{Estratégia sequencial Tx$\to$Rx completa}
\begin{algorithmic}[1]
\Statex \textbf{Etapa 1 --- executada uma vez, antes do voo}
\Require $n_1,n_2,B$, comprimento de hélice desejado $B_{helix}$,
  restrição de missão ($R_0$ ou $z_0$), peso $\gamma$
\State $\psiT^*,\beta,\rho_0,z_0 \gets$ \Call{OtimizaTx}{$\cdots$} \Comment{Algoritmo~4.1}
\State Gerar a trajetória helicoidal de Tx a partir de $\rho_0,z_0,\beta,B_{helix}$
\Statex
\Statex \textbf{Etapa 2 --- executada a cada pulso do Tx, durante o voo}
\Require malha de alvos, $\alpha_{\mathrm{res}}$, $(\beta,B_{helix})$ da Etapa~1
\For{cada posição $\mathrm{Tx}_k$ ao longo da hélice}
  \State Pré-computar $T_{1,k},R_{T,k}$ (independentes do Rx --- Proposição~3.1)
  \State Construir geometria de busca guiada por Brewster (\texttt{buildRxSearchGeometry})
  \State Resolver $\operatorname*{arg\,min}_{(x,y)} J$ (Eq.~\eqref{eq:etapa2}) com \texttt{fmincon}, 5 arranques
  \State $\mathrm{Rx}_k^* \gets$ solução
\EndFor
\State \Return trajetória completa $\{\mathrm{Tx}_k,\mathrm{Rx}_k^*\}$
\end{algorithmic}
\end{algorithm}

\section{Custo computacional}

A Etapa~1 é uma busca escalar unimodal, tipicamente $<50$ avaliações de
$J_{\mathrm{Tx}}$ até convergência — desprezível frente à Etapa~2, que
resolve um problema não-convexo de 2 variáveis, 5 arranques e até 200
iterações de \texttt{fmincon}, \emph{repetido para cada uma das dezenas de
posições decimadas do Tx} (\texttt{decimationFactor=150} em
\texttt{run\_plano\_de\_voo.m}). A Etapa~1 roda uma única vez, antes do
voo; seu custo é irrelevante frente ao da Etapa~2, o que reforça —
também do ponto de vista puramente computacional — que não há razão para
tentar fundir as duas etapas em uma busca conjunta de maior dimensão.

%%═══════════════════════════════════════════════════════════════════════════════
\chapter{Conclusões, Limitações e Trabalho Futuro}
%%═══════════════════════════════════════════════════════════════════════════════

\section{Resumo do método}

\begin{enumerate}
  \item A equação de radar biestático com refração é estruturalmente
  simétrica em Tx e Rx, mas o sistema atual só otimiza o Rx.
  \item A potência bistática se fatora \emph{exatamente} em um termo
  só-Tx e um termo só-Rx (Proposição~\ref{prop:separabilidade}) — a
  decomposição sequencial Tx$\to$Rx é, portanto, ótima para potência, sem
  aproximação.
  \item O único acoplamento residual (via $\dxy$) é absorvido pela
  liberdade posicional do Rx, já explorada pela Etapa~2 existente.
  \item O projeto do Tx se reduz a um problema escalar unimodal
  (Eq.~\eqref{eq:custo_tx}), resolúvel por busca de seção áurea, muito
  mais barato que o problema 2D não-convexo do Rx.
  \item Aplicado aos parâmetros reais do sistema, o método reproduz o
  ângulo já em uso ($58^\circ$ vs.\ $57{,}6^\circ$ previsto, $\gamma=0{,}3$) —
  validando a formalização sem contradizer o projeto existente, e
  fornecendo a ferramenta para recalcular $\psiT^*$ sempre que os
  parâmetros do sistema mudarem.
\end{enumerate}

\section{Limitações}

\begin{itemize}
  \item A relação $\theta_{i,1}\approx\psiT$ é uma aproximação de
  primeira ordem (mesma natureza da aproximação já quantificada em
  $\sim\!7$--$8\%$ para $\dxy$ no Cap.~6 da dissertação); o dimensionamento
  fino de $\rho_0,z_0$ para uma malha extensa de alvos deveria, em
  trabalho futuro, usar o ângulo de incidência exato via Newton-Raphson
  (\texttt{calculateRefractionPointFermat}), não a aproximação de
  \textit{look angle}.
  \item A Proposição~\ref{prop:aditividade} assume que $\RT,G_T$ variam
  pouco entre $\psiT^{(0)}$ e $\psiT^{(1)}$ frente à variação de $T_1$;
  para deslocamentos muito grandes de $\psiT$ essa aproximação deve ser
  revisada numericamente.
  \item Este documento não repete a validação por simulação de
  \textit{backprojection} (Cenário C, trabalho futuro identificado no
  Cap.~6 da dissertação); a validação da Etapa~1 aqui apresentada é
  analítica e numérica sobre as fórmulas já validadas, não sobre uma
  simulação eletromagnética completa.
\end{itemize}

\section{Trabalho futuro}

\begin{itemize}
  \item Implementar \texttt{optimizeTxGeometry.m} seguindo o
  Algoritmo~4.1 e integrá-lo a \texttt{run\_plano\_de\_voo.m} como um
  passo anterior ao laço principal.
  \item Estender a Etapa~1 para otimizar também o número de voltas e a
  velocidade angular da hélice, mantendo $B_{helix}$ como restrição de
  tempo de integração coerente (não tratado aqui).
  \item Quantificar, por simulação de \textit{backprojection}
  (análogo ao Cenário~A do Cap.~6 da dissertação), o ganho de SNR real
  obtido pela Etapa~1 em um caso completo, fechando o ciclo de validação
  já estabelecido para o Rx.
\end{itemize}

\end{document}
